% MAS1004 Assignment 2: long report.
%
% You fill this in with your agent, in your repository. AGENTS.md tells the
% agent how. Red text in the PDF is what still has to be replaced.
%
% - Pictures go in figures/. Copy them from results/, which is not in git.
% - Outputs go in outputs/: copy experiments.csv, LOG.md and the outputs of
%   check_submission.py and curate.py count here. \outputfile{outputs/...}
%   prints a file exactly as it is. Every number in this report has to come
%   from such an output.
%
% To make the PDF: compress this folder into a zip, and on Overleaf choose
% New Project, Upload Project. Set Menu, Compiler to XeLaTeX, and Recompile.

\documentclass[11pt]{article}
\usepackage{mas1004}

\reporttitle{Assignment 2: Long report}
\studentname{오준명}
\studentid{20261227}

\begin{document}
\makeheader

Leaderboard nickname: \texttt{ogmti}. Last releases: \texttt{track-a-1}
(``starter recipe'') and \texttt{track-b-1} (``starter model''), both made on
10 November.

% ---------------------------------------------------------------------------
\section*{Track A}

\begin{center}\footnotesize
\begin{tabular}{lp{0.6\linewidth}r}
\toprule
tag & what changed from the run before & validation \\
\midrule
baseline, seed 0 & the starter recipe as given, all training photos (first full run) & 0.4919 \\
baseline, seed 1 & only the seed & 0.4909 \\
\bottomrule
\end{tabular}
\end{center}

The only Track A recipe I trained is the starter recipe, unchanged, on all
training photos, with seeds 0 and 1. Both runs were on Colab. On my own
computer \texttt{official.py} stopped with an error on the Mac GPU (mps), and
a run on the CPU would take about 189 minutes, as \texttt{official.py -{}-quick}
printed.

The two runs differ only in the seed, and their validation scores are 0.4919
and 0.4909, a difference of 0.0010. By city, the difference is larger:

\begin{center}\footnotesize
\begin{tabular}{lrrrrr}
\toprule
 & seoul & busan & tokyo & osaka & taipei \\
\midrule
seed 0 & 0.5737 & 0.3806 & 0.6174 & 0.3398 & 0.5479 \\
seed 1 & 0.5630 & 0.3986 & 0.6138 & 0.3689 & 0.5105 \\
\bottomrule
\end{tabular}
\end{center}

\fillin{What you think these results mean.}

My final recipe is \texttt{track\_a/recipe.py} as it came, with no
\texttt{selection.csv}, released as \texttt{track-a-1}. It has no leaderboard
score yet: the release is from 10 November, and the grader scores it after
midnight.

I did not make any Track A change after the baseline (PLAN Steps 6 and 7),
because time ran out. The first release was also after the mid-point check of
30 October.

% ---------------------------------------------------------------------------
\section*{Decisions you made from photos}

Before training I looked at sheets of training photos of all five cities
(\texttt{curate.py look}). What I wrote down in \texttt{LOG.md}:

\begin{quote}
각 나라는 도시 사진 속 간판 언어의 차이에 따라 쉽게 구분할 수 있지만 간판이 없는
사진도 있기 때문에 글자 외의 요인에 의한 구분이 필요하다. 같은 나라 안의 도시는
기준을 찾기 어려웠으며 도시의 거리와 직접적으로 연관되지 않은 박물관 유물이나
자동차 실내 사진 등이 있었다. 단, 부산의 경우 바다가 자주 나오는 등의 경향을 통해
기준을 잡을 수 있었다.
\end{quote}

After the baseline, \texttt{curate.py mistakes} made one sheet per city of the
48 validation photos that baseline seed 0 got most confidently wrong. It
printed how many validation photos of each city were wrong: seoul 1038 of
2435, busan 859 of 1390, tokyo 1485 of 3881, osaka 1404 of 2128, taipei 1013
of 2243. The sheet for Busan:

\picturefile{figures/mistakes_baseline_seed0_busan_00.png}

\fillin{Two things the mistakes have in common, in your own words.}

I did not drop or add any photos: there is no \texttt{track\_a/dropped.csv}
and no \texttt{selection.csv}, so there are no scores before and after to
compare.

% ---------------------------------------------------------------------------
\section*{Track B}

One model: \texttt{src/train\_b.py} as given (ResNet18 at 224 pixels, 4
epochs, the starter training photos, \texttt{IMAGENET1K\_V1} starting
weights), trained on my Mac (mps) in 1144 seconds. Its validation score is
0.4926 (seoul 0.5363, busan 0.3755, tokyo 0.6284, osaka 0.3745, taipei
0.5484). It is released as \texttt{track-b-1}. I made no further Track B
change (PLAN Step 9).

I collected no extra photos.

There is no leaderboard score yet to put next to the validation score, for
the same reason as in Track A.

\fillin{Why the validation score and the leaderboard score might disagree.}

\subsection*{Your submission, as the grader sees it}
\outputfile{outputs/check_submission.txt}

% ---------------------------------------------------------------------------
\section*{What made the difference}

I made no change in either track, so no change can be said to have helped or
hurt. What I can compare is the Track A baseline (0.4919 and 0.4909) with
the Track B starter model (0.4926): the same network and the same photos,
trained by two different scripts, end within 0.002 of each other on the
validation photos. In all three runs Tokyo has the highest score (0.61 to
0.63) and Busan and Osaka the lowest (0.34 to 0.40).

\fillin{What you would do next with one more week.}

% ---------------------------------------------------------------------------
\section*{All runs}
\outputfile{outputs/experiments.csv}

\section*{Lab notebook}
\outputfile{outputs/LOG.md}

\end{document}
